% REWRITTEN 2026-10-08 against Ben's review of the first page. It was 492 words --
% Pie's is 212, DirectKV's 219, BOOST's 278 -- in four paragraphs, and its results
% paragraph was raw millisecond figures measured on one GPU. Three things came out
% and should stay out:
%
%  * "raises peak concurrency from 8 to 48". A VACUOUS metric: admitting requests
%    says nothing about serving them, and on the substrate where we measured that
%    number the throughput was 0.53x. Quoting the admission count without the
%    throughput reads as though we were hiding the second number.
%  * "output token-identical to the unmodified engine". Table stakes. Claiming it
%    as a result signals that we think it is one.
%  * per-millisecond coefficients and per-model gains. They belong in sections 3
%    and 6; in an abstract they make a general mechanism look like one tuning
%    result for one part number.
%
% What replaces them is the thing that makes the mechanism general: the host cohort
% is sized from what the server OBSERVES, so the same system declines to use the
% second tier where it cannot pay. That is the claim the paper can make everywhere.
\begin{abstract}
On a decode node the batch is limited by KV-cache capacity rather than by
arithmetic: the KV cache exhausts device memory well before the arithmetic units are
busy, so the server turns requests away while the GPU idles. Coherent CPU--GPU
interconnects make host DRAM fast enough to hold the KV of an \emph{actively
decoding} request, but existing systems choose \emph{how much} KV sits on the host
and never \emph{whose}: they place none, or a uniform fraction of every request, or
a bandwidth-proportional share of each memory access, or whole requests in rotation
--- pausing each one while its KV is parked. \textbf{None decides which requests
live in which tier while all of them keep running.}

We present a heterogeneous KV memory that manages HBM and host DRAM as one store and
assigns each request a tier. A \emph{host-tier} request holds no device blocks while
it generates: its KV is streamed one layer at a time, just ahead of the attention
that consumes it, and both tiers are rows of a \textbf{single forward pass}, so the
weights are read once however the batch is split. An admission controller sizes the
host cohort from the step time and link bandwidth it \emph{measures at runtime},
admitting none where the interconnect cannot hide the transfer.

We implement this in vLLM and decompose what a host-tier request costs: a fixed
per-step component, a per-request component, and a knee beyond which the step is the
transfer rather than the compute. Below the knee the per-request term carries
\emph{no} bandwidth term at all, so a faster interconnect buys a higher knee rather
than a cheaper request --- which inverts how the choice of interconnect is usually
argued. On the PCIe part we have, the model predicts and we then measure a net
\textbf{loss}; on a coherent link it projects gains of $\times1.25$ to
$\times1.37$, and we check its inputs against published Grace-Hopper measurements
rather than asking the reader to take the extrapolation on trust.
\end{abstract}
